% Hints for the coding exercise(s) of ch15 (assembled into Appendix C)

\begin{solution}{exr:ch15-flocking}
Checkpoints. \code{simulate\_\allowbreak swarm} in \code{ch15\_\allowbreak potential\_\allowbreak fields.py} already carries
separation (the pairwise term of \cref{eq:ch15-interagent}) and an optional spring, and
it builds the pairwise distance matrix once per step --- keep that matrix and read
$N_i=\set{j\ne i : \norm{\pos_i-\pos_j}\le\rho_{\mathrm{a}}}$ off it as a boolean mask
instead of looping. Positions and velocities are two $(n,2)$ arrays and the adjacency an
$(n,n)$ boolean array with a false diagonal. Steps: add cohesion as an attractive
potential towards the neighbour centroid, $\tfrac12k_{\mathrm{c}}
\norm{\pos_i-\bar\pos_i}^2$, so that it enters the same force sum as everything else;
apply the velocity update of \cref{eq:ch15-alignment} \emph{after} the force integration,
with a gain $\alpha$ small enough that $\alpha\abs{N_i}<1$ for every drone, which is the
only stability condition you need; log two scalars per step, the standard deviation of
the velocities and the radius of the position cloud; release twenty drones with random
positions and velocities. Expected behaviour: on a graph that starts connected the
velocity spread decays roughly geometrically to zero while the position spread settles at
a constant and the flock translates rigidly, at a rate set by the algebraic connectivity
--- \cref{ch:ch23} in disguise. For the split, place two tight clusters more than
$\rho_{\mathrm{a}}$ apart with velocities pointing away from each other: two components,
two limits, and a gap that only grows. Keep separation switched on and watch the minimum
pairwise distance, which stays above $1.0$ in the chapter's swarm run and drops to $0$
without it. Classic mistakes: an alignment gain not scaled by the neighbour count, so a
drone with eight neighbours overshoots; recomputing neighbourhoods from stale positions;
reading convergence off the mean velocity, which \cref{eq:ch15-alignment} conserves and
which therefore says nothing. Done when one figure shows the velocity spread going to
zero on a connected graph and holding two levels on a disconnected one, and you can name
the eigenvalue that sets the rate.
\end{solution}

\begin{solution}{exr:ch15-coding}
Checkpoints; read \code{ch15\_\allowbreak potential\_\allowbreak fields.py} only when stuck. The functions to
compare against afterwards are \code{total\_\allowbreak force} (\cref{alg:ch15-force}),
\code{simulate\_\allowbreak many} (\cref{alg:ch15-descend}), \code{numerical\_\allowbreak force},
\code{stiffness}, \code{basin\_\allowbreak experiment} and \code{follow\_\allowbreak waypoints}. One
\code{ApfParams} record holds $k_{\mathrm{att}},k_{\mathrm{rep}},\rho_0,\dt,v_{\max}$ and
the goal tolerance; obstacles are objects with a single \code{distance(p)} method, so a
segment is a drop-in replacement for a disc. Steps: write the potentials first and the
forces second, then assert that your force matches a central difference of your potential
to about $10^{-6}$ at a hundred random points --- this one test catches most sign and
chain-rule errors; add the clipped Euler step, a step budget and a stuck detector; build
the four failure scenes; sweep the basin; add the \astar layer last. Expected numbers: the
local minimum traps the drone at $x=2.487$ after $354$ steps, exactly the root of $F_x$
on the axis; GNRON stops $0.568$~m short and the exponent $n=2$ of
\cref{eq:ch15-gnron-f} arrives in $425$ steps; in the passage the stiffness of
\cref{eq:ch15-stiffness} is $667.7$ for a $0.6$~m gap and $81.0$ for a $1.0$~m gap, so
$\dt=0.05$ gives $\dt\lambda_{\max}\approx33$ and a crash while $\dt=0.01$ merely
oscillates --- the explaining number is $\dt\lambda_{\max}$, not the gap. On the
$21\times21$ grid of starts of \cref{sec:ch15-basin}, $59.4\,\%$ arrive plain and
$87.5\,\%$ with random-walk escapes, at about $40\,\%$ more steps; the waypoint layer
should beat both, and the interesting case is a waypoint inside $\rho_0$ of a wall.
Classic mistakes: an unclamped repulsion, so one step crosses the obstacle; using the
centre distance instead of the clearance of \cref{eq:ch15-clearance}; calling every
slowdown a local minimum. Done when each failure has a figure and its number, and your
arrival fractions land within a few points of the chapter's --- the Week~7 milestone
(\cref{ch:appA}).
\end{solution}
